%! Sine and Cosine Function Graphs — Unit 6, Lesson 6.5 — Group Activity
\documentclass[10pt]{article}
\usepackage{precalculus-article}
\usepackage{precalculus-boxes}
\newcommand{\TallMath}[1]{$\displaystyle #1\rule[-1.4em]{0pt}{3.2em}$}

\begin{document}

\pageheader{Unit 6, Lesson 6.5}{Group Activity}
\namepartnerperiod

\vspace{0.1in}

% Tier R
\begin{tcolorbox}[colback=white, colframe=black!40,
    title=\textbf{Tier R --- Remediate},
    fonttitle=\bfseries, arc=2mm, left=3mm, right=3mm, top=2mm, bottom=2mm]

The table below gives the values of $\sin\theta$ at key angles. Use it to answer the
questions.

\vspace{6pt}
\renewcommand{\arraystretch}{1.5}
\begin{tabular}{@{}|c|c|c|c|c|c|c|c|c|@{}}
    \hline
    $\theta$ & $0$ & $\dfrac{\pi}{6}$ & $\dfrac{\pi}{4}$ & $\dfrac{\pi}{3}$ &
               $\dfrac{\pi}{2}$ & $\pi$ & $\dfrac{3\pi}{2}$ & $2\pi$ \\\hline
    $\sin\theta$ & $0$ & $\dfrac{1}{2}$ & $\dfrac{\sqrt{2}}{2}$ & $\dfrac{\sqrt{3}}{2}$ &
                   $1$ & $0$ & $-1$ & $0$ \\\hline
\end{tabular}

\begin{enumerate}[leftmargin=1.5em, itemsep=10pt]
    \item At what values of $\theta$ in $[0, 2\pi]$ is $\sin\theta = 0$?
          These are called the \textbf{zeros} of the function.

          \blank{8cm}

    \item At what value of $\theta$ in $[0, 2\pi]$ does $\sin\theta$ reach its
          \textbf{maximum}? What is that maximum value?

          \blank{8cm}

    \item On what interval from the table does $\sin\theta$ go from $0$ to $1$? Is the
          function \textbf{increasing} or \textbf{decreasing} on that interval?

          \blank{8cm}

    \item One complete cycle of $\sin\theta$ runs from $\theta=0$ to $\theta=\blank{2cm}$.
          This length is the \textbf{period}.
\end{enumerate}
\end{tcolorbox}

\vspace{0.1in}

% Tier A
\begin{tcolorbox}[colback=white, colframe=black!40,
    title=\textbf{Tier A --- Approaching Proficiency},
    fonttitle=\bfseries, arc=2mm, left=3mm, right=3mm, top=2mm, bottom=2mm]

The table of $\cos\theta$ values below has some cells missing. Fill them in using the unit
circle or symmetry, then answer the questions.

\vspace{6pt}
\renewcommand{\arraystretch}{1.5}
\begin{tabular}{@{}|c|c|c|c|c|c|c|c|c|@{}}
    \hline
    $\theta$ & $0$ & $\dfrac{\pi}{6}$ & $\dfrac{\pi}{4}$ & $\dfrac{\pi}{3}$ &
               $\dfrac{\pi}{2}$ & $\pi$ & $\dfrac{3\pi}{2}$ & $2\pi$ \\\hline
    $\cos\theta$ & \blank{1.4cm} & $\dfrac{\sqrt{3}}{2}$ & \blank{1.4cm} &
                   $\dfrac{1}{2}$ & $0$ & \blank{1.4cm} & $0$ & \blank{1.4cm} \\\hline
\end{tabular}

\begin{enumerate}[leftmargin=1.5em, itemsep=10pt]
    \item On what interval(s) in $[0, 2\pi]$ is $\cos\theta > 0$?

          \blank{8cm}

    \item On what interval(s) in $[0, 2\pi]$ is $\cos\theta < 0$?

          \blank{8cm}

    \item A student notices that the cosine graph looks exactly like the sine graph, just
          shifted horizontally. By how much, and in which direction, would you shift the
          sine graph to get the cosine graph?

          \blank{8cm}

          Complete the equation: $\cos\theta = \sin\!\bigl(\theta + \blank{2cm}\bigr)$.
\end{enumerate}
\end{tcolorbox}

\vspace{0.1in}

% Tier E
\begin{tcolorbox}[colback=white, colframe=black!40,
    title=\textbf{Tier E --- Extension},
    fonttitle=\bfseries, arc=2mm, left=3mm, right=3mm, top=2mm, bottom=2mm]

Use the table values from Tier R to estimate the \textbf{average rate of change} of
$f(\theta) = \sin\theta$ on each interval.

\vspace{6pt}
Recall: average rate of change on $[a, b] = \dfrac{f(b) - f(a)}{b - a}$.

\vspace{8pt}
\renewcommand{\arraystretch}{2}
\begin{tabular}{@{}|l|c|@{}}
    \hline
    \textbf{Interval} & \textbf{Avg.\ rate of change} \\\hline
    $\left[0,\;\dfrac{\pi}{4}\right]$ & \blank{4cm} \\\hline
    $\left[\dfrac{\pi}{4},\;\dfrac{\pi}{2}\right]$ & \blank{4cm} \\\hline
    $\left[\dfrac{\pi}{2},\;\dfrac{3\pi}{4}\right]$ (use $\sin\tfrac{3\pi}{4}=\tfrac{\sqrt{2}}{2}$)
        & \blank{4cm} \\\hline
\end{tabular}

\begin{enumerate}[leftmargin=1.5em, itemsep=10pt, start=1]
    \item How do the three rates of change compare in size?

          \writelines{2}

    \item Near $\theta = 0$ the graph is \blank{3cm} (steep/flat), and near $\theta =
          \dfrac{\pi}{2}$ the graph is \blank{3cm} (steep/flat). Explain why, using the
          unit-circle point's motion.

          \writelines{3}
\end{enumerate}
\end{tcolorbox}

\end{document}
